\documentclass[11pt,letterpaper]{article}

\usepackage[margin=1in]{geometry}
\usepackage[T1]{fontenc}
\usepackage{palatino}
\usepackage{microtype}
\usepackage{booktabs}
\usepackage{enumitem}
\usepackage{titlesec}
\usepackage{parskip}
\usepackage{hyperref}
\usepackage{xcolor}

\hypersetup{
  colorlinks=true,
  linkcolor=blue!60!black,
  urlcolor=blue!60!black,
}

\titleformat{\section}{\large\bfseries}{\thesection.}{0.5em}{}
\titleformat{\subsection}{\normalsize\bfseries}{\thesubsection}{0.5em}{}
\titleformat{\subsubsection}{\normalsize\itshape}{\thesubsubsection}{0.5em}{}

\title{\textbf{prfaq Plugin: Architecture Overview}}
\author{Generated from codebase analysis}
\date{August 2026}

\begin{document}
\maketitle

\section{What It Is}

\texttt{prfaq} is a Claude Code plugin (v1.7.2) that implements Amazon's Working
Backwards PR/FAQ process. It generates, reviews, stress-tests, and iterates on
product discovery documents inside the terminal. The output is a professionally
typeset LaTeX PDF (or Word \texttt{.docx} via pandoc).

The plugin contains no runtime code. It is composed entirely of prompt
engineering artifacts---skill definitions, agent system prompts, reference
guides, LaTeX templates, and shell scripts. Claude Code is the runtime; the
plugin provides domain expertise and workflow structure.

\section{Project Structure}

Everything the plugin ships lives under \texttt{plugin/}, which is the plugin
root. The layout serves Claude Code's \texttt{git-subdir} marketplace source
(\texttt{"path": "plugin"}): once the entry is repointed, an install fetches
that subtree and nothing else. \texttt{\$\{CLAUDE\_PLUGIN\_ROOT\}} resolves to
it, and the paths below are relative to it. Six component categories:

\begin{description}[style=nextline, leftmargin=2em]
  \item[\texttt{.claude-plugin/plugin.json}]
    Plugin manifest declaring name, version, and description.

  \item[\texttt{skills/prfaq/}]
    One skill (\texttt{SKILL.md}) orchestrating the six-phase generation
    workflow, plus ten reference guides encoding domain knowledge.

  \item[\texttt{agents/}]
    Eight specialized subagents, each with a narrow role and specific reference
    guide dependencies.

  \item[\texttt{commands/}]
    Fourteen slash commands for post-generation operations (review, meeting,
    feedback, export, vote, etc.).

  \item[\texttt{assets/}]
    Two LaTeX templates: one for PR/FAQ documents, one for externalized press
    releases. Also includes a styled Word reference document for
    \texttt{.docx} export.

  \item[\texttt{scripts/}]
    Five shell scripts handling PDF compilation, Word export, LaTeX
    preprocessing for pandoc, reference document generation, and project
    permission grants.
\end{description}

The repo root holds development material a consumer never receives---the
dogfood \texttt{prfaq.tex}, \texttt{docs/}, \texttt{research/},
\texttt{meetings/}, \texttt{tests/}, and this repo's own release scripts in a
second, unshipped \texttt{scripts/}.

\section{The Skill}

The single skill (\texttt{/prfaq}) is a 280-line prompt that orchestrates
document generation across six phases:

\begin{enumerate}[nosep]
  \item \textbf{Phase 0: Research Discovery} --- Scans \texttt{./research/} for
    primary data, invokes the researcher agent, builds a \texttt{.bib} file.
  \item \textbf{Phase 1: Discovery} --- Interactive questioning: stage, customer,
    problem, solution, differentiation, market, risks.
  \item \textbf{Phase 2: Draft PR} --- Generates the eight press release sections
    from the LaTeX template and \texttt{pr-structure.md}.
  \item \textbf{Phase 3: Draft FAQ} --- External FAQs, internal FAQs (by risk
    category), four-risks assessment, feature appendix, then peer review.
  \item \textbf{Phase 4: Compile} --- Runs \texttt{compile\_prfaq.sh}, fixes
    overfull hbox warnings, iterates until clean.
  \item \textbf{Phase 5: Review} --- Evaluates the document against seven review
    criteria and identifies the weakest sections.
\end{enumerate}

The skill also handles revise mode: if \texttt{prfaq.tex} already exists, it
offers surgical edits instead of regenerating from scratch.

\section{The 13 Commands}

Each command is a focused, single-purpose operation that typically invokes one
agent and handles its output:

\begin{table}[h]
\centering
\small
\begin{tabular}{@{}lll@{}}
\toprule
\textbf{Command} & \textbf{Purpose} & \textbf{Agent Used} \\
\midrule
\texttt{/prfaq:import}       & Parse existing doc, launch workflow  & researcher \\
\texttt{/prfaq:badge}        & Stage-colored README badge           & (none) \\
\texttt{/prfaq:review}       & Standalone peer review               & peer-reviewer \\
\texttt{/prfaq:meeting}      & Interactive 4-persona review         & 4 meeting agents \\
\texttt{/prfaq:meeting-hive} & Autonomous consensus meeting         & 4 meeting agents \\
\texttt{/prfaq:meeting-listen} & Voiced playback of meeting         & (none, uses TTS) \\
\texttt{/prfaq:feedback}     & Surgical redraft from directives     & feedback \\
\texttt{/prfaq:research}     & Evidence search                      & researcher \\
\texttt{/prfaq:streamline}   & Remove redundancy and weasel words   & streamliner \\
\texttt{/prfaq:vote}         & Three-gate go/no-go decision         & (inline) \\
\texttt{/prfaq:export}       & Convert to \texttt{.docx} via pandoc & (none) \\
\texttt{/prfaq:externalize}  & Customer-facing press release        & (inline) \\
\texttt{/prfaq:feedback-to-us} & Anonymous satisfaction rating      & (none) \\
\bottomrule
\end{tabular}
\end{table}

The intended workflow chains these sequentially: generate $\to$ badge $\to$
review $\to$ meeting $\to$ feedback (iterate) $\to$ streamline $\to$ vote $\to$
export $\to$ externalize.

\section{The Eight Agents}

\begin{table}[h]
\centering
\small
\begin{tabular}{@{}lll@{}}
\toprule
\textbf{Agent} & \textbf{Role} & \textbf{Key Guides Loaded} \\
\midrule
peer-reviewer     & Adversarial review (Kahneman framework)     & All 10 guides \\
researcher        & Evidence search (local, quarry, web)         & None \\
feedback          & Cascading redraft engine                     & Context-dependent \\
streamliner       & Scalpel editor for prose tightening          & precise-writing \\
meeting-engineer  & Wei: feasibility risk                        & principal-engineer, four-risks \\
meeting-customer  & Priya: value risk, customer reality          & ux-bar-raiser, four-risks \\
meeting-executive & Alex: strategic fit, devil's advocate        & unit-economics, four-risks \\
meeting-builder   & Dana: ambition risk, cost of inaction        & four-risks, common-mistakes \\
\bottomrule
\end{tabular}
\end{table}

The four meeting personas embody competing priorities. Wei and Alex pull toward
caution; Dana pulls toward ambition; Priya grounds both sides in customer
reality. In interactive mode (\texttt{/prfaq:meeting}), the user resolves
disagreements. In hive mode (\texttt{/prfaq:meeting-hive}), a door-weighted
consensus mechanism resolves them autonomously.

\section{The Ten Reference Guides}

Reference guides are the shared knowledge layer. Each is a standalone document
encoding a framework, evaluation criteria, and stage-calibrated expectations.
Agents load only the guides they need.

Every guide follows the same structure: framework definition, concrete
signals and criteria, what to check in the PR/FAQ, and a stage calibration table
that adjusts expectations for hypothesis, validated, and growth stages.

\subsection{Document Structure Guides}

\begin{description}[style=nextline, leftmargin=2em]
  \item[\texttt{pr-structure.md}]
    Section-by-section blueprint for the mock press release. Defines the eight
    sections (headline through call to action) with strong/weak examples and
    common failures. Establishes that the customer quote is fictional at all
    stages by design.

  \item[\texttt{faq-structure.md}]
    Organization of external FAQs (5 required customer-facing questions) and
    internal FAQs (organized by value/market, technical, and business categories).
    Defines the LaTeX environments (\texttt{faqpair}, \texttt{featureitem}) and
    cross-referencing system (\texttt{\textbackslash faqref\{\}},
    \texttt{\textbackslash featureref\{\}}).
\end{description}

\subsection{Framework Guides}

\begin{description}[style=nextline, leftmargin=2em]
  \item[\texttt{four-risks.md}]
    Cagan's four risks (value, usability, feasibility, viability) mapped to
    PR/FAQ sections. Includes low/high risk signals, seven review criteria, and
    eight named decision outcomes (Go, Not Differentiated, TAM Too Small,
    Investment Too Risky, Not Feasible, Not Viable, Deprioritized, Iterate).

  \item[\texttt{common-mistakes.md}]
    Seven anti-patterns: skills-forward thinking, confusing speed with velocity,
    selling instead of truth-seeking, discounting competition, vague customer
    definition, great product wrong problem, and false precision in timeline
    estimates. Each has detection signals and fixes.

  \item[\texttt{decision-quality.md}]
    Twelve cognitive biases from Kahneman, Lovallo, and Sibony (HBR 2011)
    adapted to PR/FAQ review: self-interest, affect heuristic, groupthink,
    saliency, confirmation, availability, anchoring, halo effect, sunk cost,
    overconfidence, disaster neglect, and loss aversion. Each includes
    PR/FAQ-specific detection signals.
\end{description}

\subsection{Meeting Guide}

\begin{description}[style=nextline, leftmargin=2em]
  \item[\texttt{meeting-guide.md}]
    Defines the four personas, the meeting flow (pre-meeting scan $\to$ agenda
    $\to$ per-hot-spot debate $\to$ summary), synthesis guidelines for turning
    parallel agent outputs into narrative, and hive mode consensus rules.
    Hot spots are classified as one-way doors (irreversible) or two-way doors
    (reversible), with door-weighted voting: one-way doors give Wei and Alex
    veto power; two-way door ties resolve in favor of action.
\end{description}

\subsection{Risk Lens Guides}

Three guides provide deep expertise for specific risk dimensions, each used
primarily by one meeting persona:

\begin{description}[style=nextline, leftmargin=2em]
  \item[\texttt{principal-engineer.md} (Wei)]
    Feasibility risk. Covers architecture trade-offs (monolith vs.\ services,
    build vs.\ buy, consistency vs.\ availability), operational complexity at
    scale, data model decisions as the most expensive to reverse, and dependency
    risk taxonomy. Core question: ``What is the hardest unsolved technical
    problem, and does the PR/FAQ acknowledge it?''

  \item[\texttt{unit-economics.md} (Alex)]
    Viability risk. Defines CAC, LTV, LTV:CAC ratio, payback period, gross
    margin, and contribution margin. Distinguishes when unit economics are
    primary (SaaS, marketplace) from when they are secondary (internal tools,
    open source, strategic moats). Includes concrete benchmarks (LTV:CAC above
    3:1, payback under 12 months, gross margin above 60\%).

  \item[\texttt{ux-bar-raiser.md} (Priya)]
    Usability risk. Evaluates the full customer journey in five stages: awareness,
    evaluation, onboarding, first value, habit formation. Covers cognitive load
    sources, mental model alignment, the three-step Getting Started test,
    accessibility minimums, and error recovery principles.
\end{description}

\subsection{Writing Quality Guide}

\begin{description}[style=nextline, leftmargin=2em]
  \item[\texttt{precise-writing.md}]
    Editorial standard for the streamliner agent. Adapted from Vervago's
    Precision Q+A and Amazon's ``Write Like an Amazonian.'' Defines what to cut
    (redundancy, weasel words, hollow adjectives, hedge stacking, inflated
    phrases, throat-clearing sentences) and what to preserve (citations, customer
    quotes, risk assessments, numbers, structural elements). Target: 10--20\%
    shorter with zero weasel words.
\end{description}

\section{Stage Calibration}

Every document declares its stage via \texttt{\textbackslash prfaqstage\{value\}}:

\begin{description}[nosep, leftmargin=2em]
  \item[hypothesis] Pre-validation. Soft evidence acceptable. Focus on whether
    the author has identified the riskiest assumption and proposed a test.
  \item[validated] Post-interviews or post-testing. Quantitative evidence and
    specific metrics expected. Primary customer data required in FAQs.
  \item[growth] Post-launch with real users. All claims data-backed. Usage
    analytics, churn rates, and unit economics from actual operations.
\end{description}

Every reference guide, every agent, and every meeting persona adjusts its
behavior based on the document's stage. The stage calibration tables at the end
of each guide define these adjustments. This means the same guide produces three
levels of rigor without branching logic in the agent prompts.

\section{Build System}

The quality gate is compilation of every tracked LaTeX file:

\begin{enumerate}[nosep]
  \item \texttt{prfaq.tex} --- the plugin's own PR/FAQ (dogfood document)
  \item \texttt{press-release-v1.0.0.tex} --- the externalized v1.0.0 release
  \item \texttt{plugin/assets/prfaq-template.tex} --- the template used to
    generate new documents
  \item \texttt{plugin/assets/press-release-template.tex} --- the template
    \texttt{/prfaq:externalize} builds from
  \item \texttt{docs/prfaq-overview.tex} --- this document
\end{enumerate}

All must compile to valid PDFs before every commit. \texttt{make prfaq} runs
the list; \texttt{compile\_prfaq.sh} compiles one file (\texttt{pdflatex}
$\to$ \texttt{biber} $\to$ \texttt{pdflatex} $\times$ 2). The \texttt{.docx}
export pipeline preprocesses custom LaTeX environments into standard LaTeX,
then hands off to pandoc.

\section{Key Design Decisions}

\begin{enumerate}
  \item \textbf{Agents as domain experts.} Each of the eight agents has a narrow
    role with specific reference guides. The meeting agents embody named personas
    with competing risk lenses from Cagan's framework.

  \item \textbf{Reference guides as encoded knowledge.} Domain expertise is
    defined once in standalone guides and loaded by multiple agents. No
    duplication between skill prompts and agent prompts.

  \item \textbf{Stage-aware everything.} A single \texttt{\textbackslash prfaqstage\{\}}
    declaration calibrates every agent, reviewer, and meeting persona. Hypothesis
    documents get softer scrutiny than growth documents.

  \item \textbf{LaTeX as the artifact format.} Custom environments
    (\texttt{faqpair}, \texttt{featureitem}), cross-references
    (\texttt{\textbackslash faqref\{\}}, \texttt{\textbackslash featureref\{\}}),
    and bibliography support via biblatex/biber produce a professional,
    cross-referenced document.

  \item \textbf{No runtime code.} The plugin is pure prompt engineering. Shell
    scripts handle compilation and export. Claude Code provides all execution
    capability.
\end{enumerate}

\end{document}
